\chapter*{Abstract}
\thispagestyle{MyStyle}

\begin{spacing}{1.15}
This internship, carried out at ninety.\ Startup House (Sfax), builds a
system that prepares EU Carbon Border Adjustment Mechanism (CBAM)
declarations for steel and aluminium imports inside Odoo~17. The design
separates two concerns. A deterministic, unit-tested engine computes every
compliance figure, while an agentic tool-use loop, driven by a locally
hosted open-weight language model, decides which extraction, lookup and
validation tools to call when reading supplier documents. Three safeguards
protect the numeric output from model hallucination: facts already known
to the system override the model's arguments, tool results are remembered
across a job, and a measured emissions value is only written with a
verbatim quote from the source document. The system runs on two machines
linked by a private VPN. On 50 synthetic supplier documents, accuracy on
the hardest field rose from 7\% with the first model tested to 81\% (95\%
CI 66--91\%) with the final configuration. The evaluation also showed that
a standard tool-calling benchmark wrongly ruled out a usable model, so
model choice must be validated against the real execution path.

\smallskip
\noindent\textbf{Keywords:} agentic AI, large language models, tool use,
CBAM, regulatory compliance, Odoo, hallucination mitigation.

\section*{Résumé}

Ce stage, réalisé chez ninety.\ Startup House (Sfax), développe un système
qui prépare les déclarations du Mécanisme d'ajustement carbone aux
frontières de l'UE (MACF/CBAM) pour les importations d'acier et
d'aluminium dans Odoo~17. La conception sépare deux préoccupations. Un
moteur de calcul déterministe et testé produit chaque valeur de
conformité, tandis qu'une boucle agentique, pilotée par un modèle de
langage open-weight hébergé localement, choisit les outils d'extraction,
de recherche et de validation à appeler pour lire les documents des
fournisseurs. Trois garde-fous protègent le résultat numérique contre les
hallucinations du modèle : les faits déjà connus remplacent les arguments
du modèle, les résultats d'outils sont mémorisés au cours d'une tâche, et
une émission mesurée n'est enregistrée qu'accompagnée d'une citation
exacte du document source. Le système fonctionne sur deux machines
reliées par un VPN privé. Sur 50 documents synthétiques, la précision sur
le champ le plus difficile passe de 7\,\% avec le premier modèle testé à
81\,\% (IC à 95\,\% : 66--91\,\%) avec la configuration finale.
L'évaluation montre aussi qu'un benchmark standard d'appel d'outils avait
écarté à tort un modèle utilisable : le choix du modèle doit être validé
sur le chemin d'exécution réel.

\smallskip
\noindent\textbf{Mots-clés :} IA agentique, grands modèles de langage,
utilisation d'outils, MACF, conformité réglementaire, Odoo, atténuation
des hallucinations.
\end{spacing}
